\documentclass[11pt,professionalfont,aspectratio=43]{beamer}

% Assets
\usepackage[czech]{babel}				% Jazyk
\usepackage[a-2u]{pdfx}				    % Kopírování z pdfka
\usepackage{tikz}						% Schémata automatů
\usetikzlibrary{shadows, shadows.blur, shapes.geometric, arrows, positioning, overlay-beamer-styles, patterns, shadings, shapes.callouts}
% \usetikzlibrary{arrows}
\usepackage{pgfplots}
% \pgfplotsset{compat=1.15}
% \usepackage{mathrsfs}
\usepackage[export]{adjustbox}

\usepackage[T1]{fontenc}
\usefonttheme[onlymath]{serif}
% \usepackage{newtxmath}
\usepackage{pxfonts}
\usepackage{slantsc}
% \usepackage{varwidth}

\usepackage[linesnumbered,ruled,vlined]{algorithm2e}                % Pseudokód algoritmů
\usepackage[utf8]{inputenc}
% \usepackage{textcomp}
\usepackage{hyperref}
\usepackage{fancybox}

%%% FONTS
% \usefonttheme{serif}
% \usepackage{lmodern}
\usepackage{helvet}

\hypersetup{%
    pdfencoding=auto,
    pdfauthor={\insertauthor},
    pdftitle={\insertsubtitle}
}
\usepackage{float}
\usepackage{csquotes}					% české uvozovky
\usepackage{enumerate}					% enumerate environment
\usepackage{indentfirst}
\usepackage{tcolorbox}                  % Fancy boxíky
\usepackage{mathtools}
\usepackage{pifont}
\usepackage{cancel}
\usepackage{subcaption}
\usepackage{soul}
\usepackage{makecell}                   % Zalomení v buňce tabulky
\usepackage{xcolor}
\usepackage{graphicx}
\usepackage{tabularx}
\usepackage{amsmath}
\usepackage{amsthm}
\usepackage{colortbl}                   % Barvení buněk v tabulce
\usepackage{multicol}                   % Rozdělení obsahu do sloucpů

\usepackage{listings}                   % Úryvky z kódu
\definecolor{maroon}{rgb}{0.5,0,0}
\definecolor{forestgreen}{rgb}{0.13, 0.55, 0.13}

% \lstdefinelanguage{XML}
% {
%   basicstyle=\ttfamily,
%   morestring=[s]{"}{"},
%   morecomment=[s]{?}{?},
%   morecomment=[s]{!--}{--},
%   commentstyle=\color{forestgreen},
%   moredelim=[s][\color{black}]{>}{<},
%   moredelim=[s][\color{red}]{\ }{=},
%   stringstyle=\color{blue},
%   identifierstyle=\color{maroon}
% }
% Colors

\def\maincolR{0} % červená složka (0–255)
\def\maincolG{83} % zelená složka
\def\maincolB{120} % modrá složka
\definecolor{maincolor}{HTML}{8B00B5}

% \pgfmathtruncatemacro{\authorfooterbgR}{1*\maincolR}
% \pgfmathtruncatemacro{\authorfooterbgG}{1*\maincolG}
% \pgfmathtruncatemacro{\authorfooterbgB}{1*\maincolB}
% \definecolor{authorfooterbg}{RGB}{\authorfooterbgR,\authorfooterbgG,\authorfooterbgB}

% \pgfmathtruncatemacro{\titlefooterbgR}{0.722*\maincolR}   
% \pgfmathtruncatemacro{\titlefooterbgG}{0.974*\maincolG}   
% \pgfmathtruncatemacro{\titlefooterbgB}{0.739*\maincolB}
% \definecolor{titlefooterbg}{RGB}{\titlefooterbgR,\titlefooterbgG,\titlefooterbgB}

% Beamer theme
\colorlet{themecolor}{maincolor}
\colorlet{authorfooterbg}{maincolor}
\definecolor{titlefooterbg}{HTML}{9471B0}
\definecolor{titlefooterfg}{HTML}{000000}
\colorlet{titlecolorfg}{maincolor}
\definecolor{titleboxbg}{HTML}{E7E7E7}
\definecolor{datefooterbg}{HTML}{E3E3E3}
\colorlet{titleframeboxcolor}{maincolor}

% Background gradient
\definecolor{gradMiddle}{RGB}{212,248,255}
\definecolor{gradEnd}{RGB}{114,231,255}

% Block theme
\definecolor{blocktitlebg}{HTML}{F0DFA9}
\definecolor{blocktitlefg}{HTML}{957300}
\definecolor{blockbodybg}{HTML}{F5F0DF}

\definecolor{alertblockbodybg}{HTML}{FFE4E4}
\definecolor{alertblocktitlebg}{HTML}{ECA4A4}
\definecolor{alertblocktitlefg}{HTML}{B11313}

\definecolor{exampleblockbodybg}{HTML}{DAF0DA}
\definecolor{exampleblocktitlebg}{HTML}{BBE0BB}
\definecolor{exampleblocktitlefg}{HTML}{257A25}

\definecolor{mathfg}{HTML}{00185E}

\colorlet{bulletbg}{maincolor}
% \definecolor{bulletbg}{HTML}{}

% Beamer theme
\usetheme{Madrid}
\usecolortheme{seahorse}
% \useoutertheme[subsection=false]{miniframes}

\usecolortheme[named=themecolor]{structure}
\setbeamertemplate{frame numbering}[fraction]
\setbeamertemplate{navigation symbols}{}
\setbeamercovered{invisible}

% Header/footer
\setbeamercolor{author in head/foot}{bg=authorfooterbg,fg=white}
\setbeamercolor{title in head/foot}{bg=titlefooterbg,fg=titlefooterfg}
\setbeamercolor{date in head/foot}{bg=datefooterbg,fg=black}
% \setbeamercolor{title}{bg=datefooterbg}
% \setbeamercolor{title}{fg=titlecolorfg}
\setbeamerfont{title}{series=\bfseries}
\setbeamercolor{frametitle}{bg=datefooterbg}
% \setbeamercolor{section in head/foot}{bg=red,fg=cyan}
% \setbeamercolor{subsection in head/foot}{bg=blue,fg=yellow}

% Table of contents
\setbeamercolor{section in toc}{fg=black}
\setbeamercolor{subsection in toc}{fg=red}
\setbeamercolor{section number projected}{fg=black}

% Blocks
\setbeamertemplate{blocks}[rounded][shadow=true]

\setbeamercolor{block title}{bg=blocktitlebg, fg=blocktitlefg}
\setbeamercolor{block body}{parent=normal text,use=block title,bg=blockbodybg}

\setbeamercolor{block body alerted}{bg=alertblockbodybg}
\setbeamercolor{block title alerted}{bg=alertblocktitlebg,fg=alertblocktitlefg}

\setbeamercolor{block body example}{bg=exampleblockbodybg}
\setbeamercolor{block title example}{bg=exampleblocktitlebg,fg=exampleblocktitlefg}

% Math
\setbeamercolor{math text}{fg=mathfg}

% Itemize a enumerate
\setbeamertemplate{itemize items}[circle]
\setbeamertemplate{enumerate items}[circle]
\setbeamercolor{itemize item}{fg=bulletbg,bg=white}
\setbeamercolor{itemize subitem}{fg=bulletbg,bg=white}
\setbeamercolor{itemize subsubitem}{fg=bulletbg,bg=white}

\setbeamertemplate{enumerate items}[circle]
\setbeamercolor{item projected}{bg=bulletbg}

% Titlepage
\defbeamertemplate*{title page}{customized}[1][]
{
    \vbox{}%
    \vfill%
    \begin{centering}
        \begin{tikzpicture}
            \node[thick, rectangle, rounded corners=5mm, drop shadow={shadow xshift=0mm, shadow yshift=2mm, opacity=0.2, blur shadow, shadow blur steps=10}, inner sep=5mm, fill=titleboxbg, fill opacity=1] {%
                {\usebeamerfont{title}\textcolor{titlecolorfg}{\inserttitle}\par}%
            };
        \end{tikzpicture}
        \vskip1em\par%
        \ifx\insertsubtitle\@empty%
        \else%
            {\usebeamerfont{subtitle}\insertsubtitle\par}%
        \fi%
        \vskip1em\par%
        \ifx\insertauthor\@empty%
        \else%
            {\usebeamerfont{author}\insertauthor}\\%
            \texttt{\email}
        \fi
        \vskip0.3em\par%
        % \ifx\insertinstitute\@empty%
        % \else%
        %     {\usebeamerfont{institute}\insertinstitute\par}%
        % \fi%
        \vskip1em\par%
        \ifx\inserttitlegraphic\@empty%
        \else%
            \vskip1em\par%
            \inserttitlegraphic\par%
        \fi%
        \vskip1em\par%
        \ifx\insertdate\@empty%
        \else%
            {\usebeamerfont{date}\insertdate\par}%
        \fi%
    \end{centering}%
    \vfill%
}
% Enumerate
%\setlist[enumerate]{topsep=0pt,itemsep=-1ex,partopsep=1ex,parsep=1ex,label=(\arabic*)}

\MakeOuterQuote{"}

%%% Makra pro definice, věty, tvrzení, příklady, ... (vyžaduje baliček amsthm)

\theoremstyle{plain}
\newtheorem{thm}{Věta}[section]
\newtheorem{lem}[theorem]{Lemma}
\newtheorem{prop}[theorem]{Tvrzení}
\newtheorem{cor}[theorem]{Důsledek}
\newtheorem*{prop*}{Tvrzení}

\theoremstyle{definition}
\newtheorem{define}[theorem]{Definice}
% \newtheorem{example}[theorem]{Příklad}
% \newtheorem{remark}[theorem]{Poznámka}
% \newtheorem{convention}[theorem]{Úmluva}
% \newtheorem{denoting}[theorem]{Značení}

\newenvironment{defblock}[1][]{
  \begin{block}{\ifthenelse{\equal{#1}{}}
  {Definice}
  {Definice (#1)}}
}{\end{block}}
\newenvironment{thmblock}[1][]{
  \begin{block}{\ifthenelse{\equal{#1}{}}
  {Věta}
  {Věta (#1)}}
}{\end{block}}
\newenvironment{corblock}[1][]{
  \begin{block}{\ifthenelse{\equal{#1}{}}
  {Důsledek}
  {Důsledek (#1)}}
}{\end{block}}
\newenvironment{propblock}[1][]{
  \begin{block}{\ifthenelse{\equal{#1}{}}
  {Tvrzení}
  {Tvrzení (#1)}}
}{\end{block}}
\newenvironment{exblock}[1][]{
  \begin{exampleblock}{\ifthenelse{\equal{#1}{}}
  {Příklad}
  {Příklad (#1)}}
}{\end{exampleblock}}


% Colors
\definecolor{lightblue}{HTML}{009AD4}
\definecolor{lightgreen}{HTML}{68FF00}
\definecolor{darkgreen}{HTML}{00B600}
\definecolor{darkblue}{HTML}{0265b3}
\definecolor{darkred}{HTML}{DA0707}
\definecolor{lightred}{HTML}{FF5100}
\definecolor{orange}{HTML}{FFE000}
\definecolor{codeblue}{HTML}{007eed}
\definecolor{codegreen}{rgb}{0,0.6,0}
\definecolor{codegray}{rgb}{0.5,0.5,0.5}
\definecolor{codebeige}{HTML}{D4A000}
\definecolor{codepink}{HTML}{FF0055}
\definecolor{backcolour}{rgb}{0.95,0.95,0.92}

\definecolor{decproblemtitlefg}{HTML}{957300}
\definecolor{decproblembg}{HTML}{F5F0DF}

\definecolor{timportantboxcolor}{HTML}{A97100}
\definecolor{talertboxcolor}{HTML}{B70000}
\definecolor{tnoticeboxcolor}{HTML}{159C00}
\definecolor{tinfoboxcolor}{HTML}{004DFF}

\newcommand{\markred}[1]{\textcolor{lightred}{#1}}
\newcommand{\markdarkred}[1]{\textcolor{darkred}{#1}}
\newcommand{\markgreen}[1]{\textcolor{lightgreen}{#1}}
\newcommand{\markdarkgreen}[1]{\textcolor{darkgreen}{#1}}
\newcommand{\markorange}[1]{\textcolor{orange}{#1}}
\newcommand{\markblue}[1]{\textcolor{lightblue}{#1}}
\newcommand{\markdarkblue}[1]{\textcolor{darkblue}{#1}}

% Math mode settings
% \AtBeginDocument{
%     \setlength{\abovedisplayskip}{3pt}
%     \setlength{\belowdisplayskip}{-4pt}
% }

% Inline images
\newcommand{\inlineimgscale}{1.1}

% X and check mark
\newcommand{\cmark}{\markgreen{\ding{51}}}
\newcommand{\xmark}{\markred{\ding{55}}}

% Math
\newcommand{\R}{\mathbb{R}}
\newcommand{\C}{\mathbb{C}}
\newcommand{\N}{\mathbb{N}}
\newcommand{\Z}{\mathbb{Z}}
\newcommand{\Q}{\mathbb{Q}}

\newcommand{\T}[1]{#1^\top}                     % Tranpozice matice
\newcommand{\compl}[1]{#1^\complement}
\newcommand{\mapping}[3]{#1:#2 \rightarrow #3}
\newcommand{\mapto}[3]{#1:#2 \mapsto #3}
\newcommand{\set}[1]{\left\{#1\right\}}
\newcommand{\powset}{\mathcal{P}}
\newcommand{\code}[1]{\langle#1\rangle}
\DeclareMathOperator*{\argmin}{arg\,min}
\DeclareMathOperator*{\argmax}{arg\,max}

% Statistika a AI
\newcommand{\average}[1]{\overline{#1}}
\DeclareMathOperator{\median}{med}
\DeclareMathOperator{\Var}{var}
\DeclareMathOperator{\Cov}{cov}
\DeclareMathOperator{\MSE}{MSE}
\DeclareMathOperator{\SSE}{SSE}

\DeclareMathOperator{\TP}{TP}
\DeclareMathOperator{\TN}{TN}
\DeclareMathOperator{\FP}{FP}
\DeclareMathOperator{\FN}{FN}
\DeclareMathOperator{\Accuracy}{Acc}
\DeclareMathOperator{\Precision}{Prec}
\DeclareMathOperator{\Recall}{Rec}
\DeclareMathOperator{\Fscore}{F}
\DeclareMathOperator{\Gini}{Gini}

\DeclareMathOperator{\mspan}{span}
\DeclareMathOperator{\mrank}{rank}
\DeclareMathOperator{\tg}{tg}
\DeclareMathOperator{\id}{id}
\DeclareMathOperator{\dom}{dom}
\DeclareMathOperator{\rng}{rng}
\renewcommand{\emptyset}{\varnothing}
\newcommand{\binary}[1]{\ensuremath{\left\langle #1\right\rangle}_B}

% DFS in and out lists
\DeclareMathOperator{\inlist}{in}
\DeclareMathOperator{\outlist}{out}
% \DeclareMathOperator{\lowlist}{low}
\DeclareMathOperator{\key}{key}

% A* macro
\newcommand{\Astar}{$\textcolor{black}{\textup{A}^{*}}$}

% Complexity classes
\DeclareMathOperator{\classP}{P}
\DeclareMathOperator{\classNP}{NP}
\DeclareMathOperator{\classTime}{TIME}
\DeclareMathOperator{\classNTime}{NTIME}
\DeclareMathOperator{\classSpace}{SPACE}
\DeclareMathOperator{\classNSpace}{NSPACE}

\DeclareMathOperator{\SAT}{SAT}
\DeclareMathOperator{\THREESAT}{3-SAT}
\DeclareMathOperator{\UNSAT}{UNSAT}
\DeclareMathOperator{\NZMNA}{NzMna}
\DeclareMathOperator{\KLIKA}{Klika}
\DeclareMathOperator{\ROZVRH}{Rozvrh}
\DeclareMathOperator{\SUDOKU}{Sudoku}
\newcommand{\HALT}{\ensuremath{\textsc{Halt}}}
\newcommand{\FACTOR}{\ensuremath{\textsc{Factor}}}

\DeclareMathOperator{\regex}{RegE}

%%% Barevné boxíky (notice / info / important / alert)
%
% Šířka boxu se nezadává ručně -- box se přizpůsobí šířce svého obsahu.
% Obsah se vysází do prostředí varwidth, takže krátký text zůstane na
% jednom řádku; delší text se zalomí, nejvýše však na šířku \linewidth
% (tj. \textwidth mimo sloupce a minipage). Pokud je titulek širší než
% obsah, řídí šířku boxu titulek.

\newsavebox{\tboxcontentbox}
\newsavebox{\tboxtitlebox}
\newlength{\tboxwidth}
\newlength{\tboxpadding}

% Geometrie boxu (musí odpovídat stylu tautobox níže):
% 2 * (boxrule + boxsep + left)
\setlength{\tboxpadding}{\dimexpr 0.5mm+1mm+4mm\relax}
\setlength{\tboxpadding}{2\tboxpadding}

\tcbset{tautobox/.style={boxrule=0.5mm, boxsep=1mm, left=4mm, right=4mm,
                         fonttitle=\bfseries}}

% Začátek sběru obsahu do boxu
\newcommand{\tboxcollect}{%
  \begin{lrbox}{\tboxcontentbox}%
    \begin{varwidth}{\dimexpr\linewidth-\tboxpadding\relax}%
}

% Ukončení sběru a výpočet výsledné šířky
% #1 = titulek (pro měření; prázdný u variant bez titulku)
\newcommand{\tboxmeasure}[1]{%
    \end{varwidth}%
  \end{lrbox}%
  \sbox{\tboxtitlebox}{\bfseries #1}%
  \setlength{\tboxwidth}{\wd\tboxcontentbox}%
  \ifdim\wd\tboxtitlebox>\tboxwidth
    \setlength{\tboxwidth}{\wd\tboxtitlebox}%
  \fi
  \addtolength{\tboxwidth}{\tboxpadding}%
  \ifdim\tboxwidth>\linewidth
    \setlength{\tboxwidth}{\linewidth}%
  \fi
}

% Vysázení boxu s titulkem: #1 = barva, #2 = titulek
\newcommand{\tboxrenderbox}[2]{%
  \tboxmeasure{#2}%
  \begin{center}%
    \begin{tcolorbox}[tautobox, colback=#1!15, colframe=#1,
                      title={#2}, width=\tboxwidth]
      \usebox{\tboxcontentbox}%
    \end{tcolorbox}%
  \end{center}%
}

% Vysázení boxu bez titulku: #1 = barva
\newcommand{\tboxrenderframe}[1]{%
  \tboxmeasure{}%
  \begin{center}%
    \begin{tcolorbox}[tautobox, colback=#1!15, colframe=#1,
                      width=\tboxwidth]
      \usebox{\tboxcontentbox}%
    \end{tcolorbox}%
  \end{center}%
}

% Poznámky (notice)
\newenvironment{tnoticebox}[1]{%
  \def\tboxtitletext{#1}\tboxcollect
}{%
  \tboxrenderbox{tnoticeboxcolor}{\tboxtitletext}%
}

\newenvironment{tnoticeframe}{%
  \tboxcollect
}{%
  \tboxrenderframe{tnoticeboxcolor}%
}

% Informační bloky (info)
\newenvironment{tinfobox}[1]{%
  \def\tboxtitletext{#1}\tboxcollect
}{%
  \tboxrenderbox{tinfoboxcolor}{\tboxtitletext}%
}

\newenvironment{tinfoframe}{%
  \tboxcollect
}{%
  \tboxrenderframe{tinfoboxcolor}%
}

% Důležité bloky (important)
\newenvironment{timportantbox}[1]{%
  \def\tboxtitletext{#1}\tboxcollect
}{%
  \tboxrenderbox{timportantboxcolor}{\tboxtitletext}%
}

\newenvironment{timportantframe}{%
  \tboxcollect
}{%
  \tboxrenderframe{timportantboxcolor}%
}

% Varovné bloky (alert)
\newenvironment{talertbox}[1]{%
  \def\tboxtitletext{#1}\tboxcollect
}{%
  \tboxrenderbox{talertboxcolor}{\tboxtitletext}%
}

\newenvironment{talertframe}{%
  \tboxcollect
}{%
  \tboxrenderframe{talertboxcolor}%
}

% Decision problem
\newcommand{\decproblem}[3]{
    \begin{center}
        \begin{minipage}{.9\textwidth}
            \begin{table}[!htbp]
                \centering
                \bgroup
                \def\arraystretch{1.8}
                \begin{tabularx}{\textwidth}{|rX|}
                \hline
                \rowcolor{decproblembg} \multicolumn{2}{|c|}{\textcolor{decproblemtitlefg}{\textsc{#1}}} \\ \hline
                \markdarkred{\textbf{Vstup:}} & #2                \\ \hline
                \rowcolor{decproblembg} \markdarkred{\textbf{Otázka:}} & #3               \\ \hline
                \end{tabularx}
                \egroup
            \end{table}
        \end{minipage}
    \end{center}
}
\newcommand{\bigO}{\mathcal{O}}

% TODO
\newcommand{\todo}[1]{\textcolor{red}{(\noindent TODO: #1.)}}


\newcommand{\hlcode}[1]{\colorbox{red}{#1}}

% Algorithms
\numberwithin{algocf}{section}
\renewcommand{\algorithmcfname}{Algoritmus}
\SetKwInput{KwIn}{Vstup}
\SetKwInput{KwOut}{Výstup}
\SetKwProg{Function}{Funkce}{}{end}
\setlength{\algomargin}{25pt}
\renewcommand{\thealgocf}{}     % Vypnutí číslování algoritmů

% Definice stylů pro uzly
\tikzstyle{startstop} = [rectangle, rounded corners, minimum width=1.5cm, minimum height=0.5cm,text centered, draw=black, fill=red!30]
\tikzstyle{process}   = [rectangle, minimum width=1.5cm, minimum height=0.5cm, text centered, draw=black, fill=orange!30]
\tikzstyle{decision}  = [diamond, aspect=2, text centered, draw=black, fill=green!30]
\tikzstyle{arrow}     = [thick,->,>=stealth]

\tikzstyle{vertex}    = [thick, draw, circle, minimum size=4mm, inner sep=1pt, outer sep=1pt, align=center]
\tikzstyle{vertexplain}    = [minimum size=4mm, inner sep=1pt, outer sep=1pt, align=center]
\tikzstyle{verteximg}    = [draw, circle, minimum size=4mm, inner sep=0pt, outer sep=1pt, align=center, clip]
\tikzstyle{verteximgplain}    = [minimum size=4mm, inner sep=0pt, outer sep=1pt, align=center]
\tikzstyle{task}    = [draw, rounded corners=2pt, minimum width=14mm, minimum height=8mm, inner sep=2pt, align=center]

\tikzstyle{edge}    = [thick,>=stealth]
\tikzstyle{diredge}    = [thick, ->, >=stealth]

%% Automaty
\tikzstyle{state}=[%
    draw,
    circle,
    thick,
    minimum size=10mm,
    inner sep=0pt,
    outer sep=0pt
]

\tikzstyle{acceptstate}=[%
    state,
    double,
    double distance=1pt
]

\tikzstyle{transition}=[%
    ->,
    >=stealth,
    thick
]

\newcommand<>{\state}[4][]{%
    \node#5[state,#1] (#2) at #3 {#4};
}

\newcommand<>{\acceptingstate}[4][]{%
    \node#5[acceptstate,#1] (#2) at #3 {#4};
}

\newcommand<>{\transitionarrow}[5][]{%
    \path#6 (#2) edge[transition,#1] node[#3] {#4} (#5);
}

\newcommand<>{\initialstate}[6][]{%
    \node#7[state,#1] (#2) at #3 {#4};
    \draw#7[transition] #5 -- (#2.#6);
}

%%%%%

% Frames
\newcommand{\tableofcontentsframe}{
    \begin{frame}{Obsah}
        \tableofcontents
    \end{frame}
}

% \definecolor{boxcolor}{HTML}{00820a}
\newcommand{\titleframe}[1]{%
  \begin{frame}[plain]
    \begin{tikzpicture}[remember picture, overlay]
      \node at (current page.center)
        [draw=titleframeboxcolor, line width=1pt,
         inner xsep=10mm, inner ysep=4mm, rounded corners=2mm,
         text width=0.8\textwidth, align=center,
         execute at begin node=\setlength{\baselineskip}{2.2em}]
        {\Huge\markred{#1}\normalfont};
    \end{tikzpicture}
  \end{frame}
  \addtocounter{framenumber}{-1}%
}

\newenvironment{titleframeenv}[1]{
    \begin{frame}[plain]
        \begin{tcolorbox}[colback=white,boxrule=1pt]
            \begin{center}
                \Huge\markred{#1}\normalfont
            \end{center}
        \end{tcolorbox}
}{\addtocounter{framenumber}{-1}\end{frame}}

\newcommand{\icon}[1]{
  \includegraphics[height=\fontcharht\font`\B]{#1}
}

% Obrázek, který se vždy vejde na snímek. Adjustbox jej v případě potřeby
% zmenší tak, aby nepřesáhl šířku textu ani zadanou část výšky snímku; díky
% tomu sedí sazba ve všech používaných poměrech stran (4:3, 16:9 i 16:10).
% Volitelným argumentem je maximální výška (výchozí 0.68\textheight), povinným
% cesta k souboru s obrázkem (TikZ zdroj vkládaný přes \input).
\newcommand{\obrazek}[2][0.68\textheight]{%
    \begin{figure}
        \centering
        \begin{adjustbox}{max width=\linewidth, max totalheight=#1, center}
            \input{#2}
        \end{adjustbox}
    \end{figure}%
}

% Totéž pro rastrové či vektorové obrázky vkládané přes \includegraphics.
\newcommand{\obrazekgrf}[2][0.68\textheight]{%
    \begin{figure}
        \centering
        \includegraphics[width=\linewidth, max totalheight=#1, center]{#2}
    \end{figure}%
}

% Závěrečný snímek se seznamem zdrojů. Prvním (volitelným) argumentem je název
% sekce i snímku, druhým čárkou oddělený seznam klíčů z ../assets/literatura.bib.
% Citace se sázejí podle ČSN ISO 690, v pořadí, v jakém jsou klíče uvedeny.
\newcommand{\zdrojeframe}[2][Zdroje]{%
    \section{#1}
    \nocite{#2}%
    \begin{frame}[allowframebreaks]{#1}
        \printbibliography[heading=none]
    \end{frame}
}
\input{../assets/listing.tex}

\def\presentationtitle{Formální gramatiky}

\begin{document}

    \title[\presentationtitle]{\textsc{Teoretická informatika}}
\subtitle{\presentationtitle}
\author{David Weber}
\def\email{weber3@spsejecna.cz}
\institute{SPŠE Ječná}
\date{%
    Rok 
    \pgfmathsetmacro{\firstyear}{ifthenelse(\month<9, \year-1, \year)}%
    \pgfmathsetmacro{\secondyear}{ifthenelse(\month<9, \year, \year+1)}%
    \pgfmathtruncatemacro{\secondyearlasttwo}{mod(\secondyear,100)}%
    \pgfmathprintnumber[fixed,precision=0,zerofill,1000 sep={}]{\firstyear}/%
    \pgfmathprintnumber[fixed,precision=0,zerofill,1000 sep={}]{\secondyearlasttwo}%
}
\titlegraphic{\includegraphics[width=0.25\textwidth]{../images/SPSE-Jecna_Logo.pdf}}

\begin{frame}[plain]
    \titlepage
\end{frame}

    \section{Formální gramatiky}

    \titleframe{Formální gramatiky}

    \subsection{Rekapitulace}

    \begin{frame}{Rekapitulace}
        \begin{itemize}
            \visible<2->{\item Již víme, že jazyky můžeme popsat pomocí}
            \begin{itemize}
                \visible<3->{\item \textbf{(ne)deterministického konečného automatu} a}
                \visible<4->{\item \textbf{regulárního výrazu}.}
            \end{itemize}
            \visible<4->{\item Hodilo by se nám ale vědět, jak jazyk "vzniká".}
        \end{itemize}
        \medskip
        \visible<5->{\begin{timportantframe}
            \begin{center}
                K tomu se nám budou hodit právě tzv. \textbf{formální gramatiky}.
            \end{center}
        \end{timportantframe}}
        \medskip
        \visible<6->{\begin{tnoticebox}{Hlavní myšlenka}
            \begin{center}
                Automat slova rozpoznává, gramatika slova generuje.
            \end{center}
        \end{tnoticebox}}
    \end{frame}

    \subsection{Proč gramatiky?}

    \begin{frame}{Proč gramatiky?}
        \begin{itemize}
            \visible<2->{\item Popisují, jak mohou vznikat správně utvořená slova.}
            \visible<3->{\item Umožňují popsat strukturu vět, výrazů nebo programů.}
            \visible<4->{\item Používají se například:}
            \begin{itemize}
                \visible<5->{\item v překladačích,}
                \visible<6->{\item při syntaktické analýze,}
                \visible<7->{\item při zpracování přirozeného jazyka (tzv. počítačová lingvistika).}
            \end{itemize}
        \end{itemize}
    \end{frame}

    \subsection{Definice gramatiky}

    \begin{frame}{Definice formální gramatiky}
        \visible<2->{\begin{defblock}[Formální gramatika]
            \visible<3->{\markdarkred{Formální gramatikou} nazýváme uspořádanou čtveřici $G=(V,T,P,S)$, kde}
            \begin{itemize}
                \visible<4->{\item $V$ je množina tzv. \textbf{neterminálů},}
                \visible<5->{\item $T$ je množina tzv. \textbf{terminálů},}
                \visible<6->{\item $S\in V$ je počáteční symbol a}
                \visible<7->{\item $P$ je množina \textbf{přepisovacích pravidel} tvaru $\alpha\to\omega$ (řetězec $\alpha$ se přepíše na řetězec $\omega$), přičemž $\alpha,\omega\in(V\cup T)^*$ a $\alpha$ obsahuje alespoň jeden neterminál.} 
            \end{itemize}
            \visible<8->{Množiny $V$ a $T$ jsou disjunktní, tj. $V\cap T=\emptyset$.}
        \end{defblock}}
        \medskip
        \visible<9->{\begin{talertframe}
            \begin{center}
                Není přípustné, aby se v~přepisovacím pravidle vyskytovaly na levé straně $\alpha$ pouze terminály.
            \end{center}
        \end{talertframe}}
    \end{frame}

    \subsection{Derivace a jazyk generovaný gramatikou}

    \begin{frame}{Derivace}
        \begin{itemize}
            \pause
            \item Pokud lze slovo \(\alpha\) jedním použitím pravidla z $P$ přepsat na slovo \(\beta\),
            píšeme
            \[
                \alpha \Rightarrow_G \beta.
            \]
            \pause 

            \item Pokud lze \(\alpha\) přepsat na \(\omega\) v~několika krocích, píšeme
            \[
                \alpha \Rightarrow_G^* \omega.
            \]
            tzn. $\alpha = \beta_1 \Rightarrow \beta_2 \Rightarrow \cdots \Rightarrow \beta_n = \omega$ pro nějaké řetězce $\beta_1, \beta_2, \ldots, \beta_n\in(V\cup T)^*$.
            \pause
            \item Posloupnost $\beta_1, \beta_2, \ldots, \beta_n$ nazýváme \textbf{derivací}.
        \end{itemize}
    \end{frame}

    \begin{frame}{Příklad gramatiky a derivace slova}
        \pause
        Mějme gramatiku $G=(V,T,P,S)$, kde
        \[
        V = \set{S}, \quad T = \set{a, b}\quad\text{a}\quad P = \set{S \to aS, S \to abS, S \to abb}.
        \]

        \pause
        
        Začneme počátečním symbolem a postupně přepisujeme podle pravidel.

        \[
            S \Rightarrow aS \Rightarrow abS \Rightarrow abb
        \]

        \pause

        Výsledkem je slovo složené pouze z terminálních symbolů.
    \end{frame}

    \begin{frame}{Jazyk generovaný gramatikou}
        \pause
        \begin{defblock}[Jazyk generovaný gramatikou]
            Jazyk generovaný gramatikou \(G=(V,T,P,S)\) je množina
            \[
                L(G)=\set{w\in T^* : S\Rightarrow_G^* w}.
            \]
        \end{defblock}

        \pause

        Tedy $L(G)$ obsahuje \textbf{všechny řetězce složené z terminálních symbolů}, které lze odvodit z počátečního symbolu $S$ pomocí pravidel z $P$.
    \end{frame}

    \subsection{Příklady}

    \begin{frame}{Příklad}{všechna binární slova}

        \pause
        Mějme gramatiku

        \[
            S\to 0S\mid 1S\mid \lambda.
        \]

        \pause

        Potom $L(G)=\{0,1\}^*$.

        \pause
        \medskip
        Ukázka derivace:

        \[
            S\Rightarrow 1S\Rightarrow 10S\Rightarrow 101S\Rightarrow 101.
        \]
    \end{frame}

    \begin{frame}{Příklad}{jazyk \texorpdfstring{\(0^n1^n\)}{0n1n}}
        \pause
        Mějme gramatiku

        \[
            S\to 0S1\mid \lambda.
        \]

        \pause

        Potom

        \[
            L(G)=\set{0^n1^n : n\geqslant 0}.
        \]

        \pause

        Například:

        \[
            S \Rightarrow 0S1 \Rightarrow 00S11 \Rightarrow 0011.
        \]

        Tento jazyk není regulární.
    \end{frame}

    \begin{frame}{Příklad}{jazyk \texorpdfstring{\(a^nb^m\)}{anbm}}
        \pause
        Uvažujme gramatiku
        \[
            G=(V,T,P,S),
        \]
        kde
        \[
            V=\set{S,A,B}, \qquad T=\set{a,b}
        \]
        a pravidla jsou
        \[
        \begin{aligned}
            S &\to AB,\\
            A &\to aA \mid a,\\
            B &\to bB \mid b.
        \end{aligned}
        \]

        \pause

        Tato gramatika generuje jazyk $L(G)=\set{a^n b^m : n,m\geqslant 1}$.

        \pause

        Například:
        \[
            S \Rightarrow AB \Rightarrow aAB \Rightarrow aaB
            \Rightarrow aabB \Rightarrow aabb.
        \]

        \pause

        Neterminál \(A\) generuje blok písmen \(a\), zatímco neterminál \(B\)
        generuje blok písmen \(b\).
    \end{frame}

    \subsection{Gramatiky a přirozený jazyk}

    \begin{frame}{Gramatiky a přirozený jazyk}
        \pause
        Gramatiky vznikly také jako pokus formálně popsat strukturu přirozeného jazyka.

        \pause

        Například jednoduchá anglická (formální) gramatika může obsahovat pravidla:
        \begin{align*}
            \text{sentence}&\to \text{noun phrase}\mid\text{verb phrase}\\
            \text{noun phrase}&\to \text{article}\mid\text{noun}\\
            \text{verb phrase}&\to \text{verb}\mid\text{noun phrase}\\
            \text{article}&\to \text{the}\\
            \text{noun}&\to \text{dog}\mid \text{bone},\\
            \text{verb}&\to \text{ate}
        \end{align*}

        \pause

        Tato gramatika například generuje větu:

        \[
            \text{the dog ate the bone}
        \]

        \pause

        Gramatika ale popisuje strukturu, nikoli význam.
    \end{frame}

    \begin{frame}{Příklad gramatiky pro anglickou větu}
    \centering
    \begin{tikzpicture}[
    x=0.72cm, y=0.92cm,
    >=stealth,
    arr/.style={->, thick, draw=black!75},
    ntbox/.style={draw=none, fill=orange!85, minimum width=2.15cm,
        minimum height=0.8cm, align=center, font=\scriptsize},
    topbox/.style={draw=none, fill=yellow!80, rounded corners=3pt,
        minimum width=2.4cm, minimum height=0.85cm, align=center, font=\scriptsize},
    termbox/.style={draw=none, fill=green!70, minimum width=1.65cm,
        minimum height=0.75cm, align=center, font=\scriptsize},
    rulelabel/.style={font=\tiny, fill=white, inner sep=1pt,
        text=black!75, align=center}
]

\node[topbox] (sentence) at (10.5,7.25) {sentence};

\node[ntbox] (np1) at (7.2,5.55) {noun phrase};
\node[ntbox] (vp)  at (13.8,5.55) {verb phrase};

\node[ntbox] (art1)  at (4.6,3.75) {article};
\node[ntbox] (noun1) at (7.9,3.75) {noun};
\node[ntbox] (verb)  at (12.3,3.75) {verb};
\node[ntbox] (np2)   at (15.8,3.75) {noun phrase};

\node[ntbox] (art2)  at (14.2,2.0) {article};
\node[ntbox] (noun2) at (17.6,2.0) {noun};

\node[termbox] (the1) at (4.6,0.25) {the};
\node[termbox] (dog)  at (7.9,0.25) {dog};
\node[termbox] (ate)  at (12.3,0.25) {ate};
\node[termbox] (the2) at (14.2,0.25) {the};
\node[termbox] (bone) at (17.6,0.25) {bone};

\draw[arr] (sentence) -- (np1);
\draw[arr] (sentence) -- (vp);

\draw[arr] (np1) -- (art1);
\draw[arr] (np1) -- (noun1);

\draw[arr] (vp) -- (verb);
\draw[arr] (vp) -- (np2);

\draw[arr] (np2) -- (art2);
\draw[arr] (np2) -- (noun2);

\draw[arr] (art1) -- (the1);
\draw[arr] (noun1) -- (dog);
\draw[arr] (verb) -- (ate);
\draw[arr] (art2) -- (the2);
\draw[arr] (noun2) -- (bone);

%------------------------------------------------
% Popisky pravidel
%------------------------------------------------
\node[rulelabel] at (10.5,6.45)
    {sentence \(\to\) noun phrase verb phrase};

\node[rulelabel] at (6.25,4.75)
    {noun phrase \(\to\) article noun};

\node[rulelabel] at (14.0,4.75)
    {verb phrase \(\to\) verb noun phrase};

\node[rulelabel] at (15.9,3.0)
    {noun phrase \(\to\) article noun};

% spodní popisky – zarovnané se svislými hranami
\node[rulelabel] at (4.6,1.35)
    {article \(\to\) the};

\node[rulelabel] at (7.9,1.35)
    {noun \(\to\) dog};

\node[rulelabel] at (12.3,1.35)
    {verb \(\to\) ate};

\node[rulelabel] at (14.2,1.12)
    {article \(\to\) the};

\node[rulelabel] at (17.6,1.12)
{noun \(\to\) bone};

\end{tikzpicture}

    \end{frame}

    \section{Regulární gramatiky}

    \titleframe{Regulární gramatiky}

    \begin{frame}{Regulární gramatiky}
        \pause
        \begin{defblock}[Regulární gramatika]
            Gramatika $G=(V,T,P,S)$ je \markdarkred{regulární}, pokud každé přepisovací pravidlo z $P$ má tvar
            \[A\to \omega B\quad\text{nebo}\quad A\to\omega,\]
            kde $A,B\in V$ a $\omega\in T^*$.
        \end{defblock}
        \pause
        Příklady:
        \begin{itemize}
            \pause
            \item $S\to aS\mid bS\mid \lambda$ \markdarkgreen{je} regulární gramatika.
            \pause
            \item $S\to aSb\mid \lambda$ \markdarkred{není} regulární gramatika.
            \pause
            \item $S\to aA\mid bB,\ A\to aS\mid bA\mid \lambda,\ B\to bS\mid aB\mid \lambda$ \markdarkgreen{je} regulární gramatika.
            \pause
            \item $S\to aA\mid bB,\ A\to aS\mid bA,\ B\to bS\mid aB$ \markdarkred{není} regulární gramatika.
        \end{itemize}
    \end{frame}

    \begin{frame}{Regulární gramatiky a regulární jazyky}
        \pause
        \begin{thmblock}
            Jazyk $L$ je regulární právě tehdy, když existuje regulární gramatika $G$ taková, že $L=L(G)$.
        \end{thmblock}
        \pause
        Náznak důkazu:
        \pause
        \begin{enumerate}
            \item konečný automat $\Rightarrow$ regulární gramatika.
            \item regulární gramatika $\Rightarrow$ konečný automat,\\
        \end{enumerate}
        \pause
        \begin{timportantframe}
            \begin{center}
                Tedy regulární gramatiky generují právě regulární jazyky.
            \end{center}
        \end{timportantframe}
        \pause
        \begin{tinfoframe}
            \begin{center}
                Ukážeme pouze první bod, druhý je analogický.
            \end{center}
        \end{tinfoframe}
    \end{frame}

    \begin{frame}{Konečný automat \(\Rightarrow\) regulární gramatika}
        Mějme DFA $A=(Q,\Sigma,\delta,q_0,F)$. \pause Sestrojíme gramatiku $G=(V,T,P,S)$, kde
        \[
            V=Q,\qquad T=\Sigma,\qquad S=q_0.
        \]

        \pause

        Pravidla \(P\) definujeme takto:
        \begin{align*}
            \delta(p,a)=q
            \quad &\Longrightarrow \quad
            p\to aq\\
            p\in F
            \quad &\Longrightarrow \quad
            p\to \lambda.
        \end{align*}

        \pause

        Přijímající výpočet automatu
        \[
            q_0 \xrightarrow{a_1} q_1
            \xrightarrow{a_2} \cdots
            \xrightarrow{a_n} q_n,\qquad q_n\in F
        \]
        odpovídá derivaci
        \[
            q_0 \Rightarrow a_1q_1
            \Rightarrow a_1a_2q_2
            \Rightarrow \cdots
            \Rightarrow a_1\cdots a_nq_n
            \Rightarrow a_1\cdots a_n.
        \]

        \pause

        Tedy $L(A)=L(G)$.
    \end{frame}

    \begin{frame}{Znázornění převodu automatu na regulární gramatiku}
    \centering
    \begin{tikzpicture}[scale=0.9, every node/.style={font=\small}]

    % levé popisky
    \node[anchor=west] at (-4.8,2.3) {\textbf{Konečný automat}};
    \node[anchor=west] at (-4.8,-2.6) {\textbf{Regulární gramatika}};

    %------------------------------------------------
    % Automat vpravo nahoře
    %------------------------------------------------
    \initialstate{q0}{(1.0,2.3)}{$q_0$}{(0.0,2.3)}{west}
    \state{q1}{(3.0,2.3)}{$q_1$}
    \acceptingstate{q2}{(5.0,2.3)}{$q_2$}

    \transitionarrow[loop above]{q0}{above}{$a$}{q0}
    \transitionarrow[bend left=15]{q0}{above}{$b$}{q1}

    \transitionarrow[loop above]{q1}{above}{$b$}{q1}
    \transitionarrow[bend left=15]{q1}{above}{$a$}{q2}

    \transitionarrow[loop above]{q2}{above}{$a,b$}{q2}

    %------------------------------------------------
    % Šipka převodu mezi oběma částmi
    %------------------------------------------------
    \draw[transition, line width=1.1pt] (3.0,1.25) -- (3.0,0.25);
    \node[anchor=west] at (3.15,0.75) {\small převod};

    %------------------------------------------------
    % Gramatika vpravo dole
    %------------------------------------------------
    \node[anchor=north west, align=left, text width=7.7cm] at (-1.1,-0.05) {%
        \[
            G=(V,T,P,S)
        \]
        \[
            V=\{q_0,q_1,q_2\},\qquad
            T=\{a,b\},\qquad
            S=q_0
        \]
        \[
        \begin{aligned}
            P=\{&
            q_0 \to aq_0 \mid bq_1,\\
            &
            q_1 \to bq_1 \mid aq_2,\\
            &
            q_2 \to aq_2 \mid bq_2 \mid \lambda
            \}
        \end{aligned}
        \]
    };

\end{tikzpicture}

    \end{frame}

    \begin{frame}{Jeden jazyk třemi způsoby}
        Uvažujme jazyk $L=\set{w\in\{a,b\}^* \mid w \text{ končí na } ab}$.

        \vspace{0.3cm}

        \begin{columns}[T]
            %----------------------------------------
            % Levý sloupec
            %----------------------------------------
            \begin{column}{0.52\textwidth}
                \textbf{Regulární výraz:}
                \[
                    (a\mid b)^*ab
                \]

                \vspace{0.4cm}

                \textbf{Regulární gramatika:}

                \[
                    G=(V,T,P,S),
                \]
                kde $V=\set{S,A,B}, T=\set{a,b}$ a~pravidla jsou
                \[
                \begin{aligned}
                    S &\to aA \mid bS,\\
                    A &\to aA \mid bB,\\
                    B &\to aA \mid bS \mid \lambda.
                \end{aligned}
                \]

                \vspace{0.2cm}

                \textbf{Příklad slov z jazyka:}
                \[
                    ab,\quad aab,\quad bab,\quad bbab.
                \]
            \end{column}

            %----------------------------------------
            % Pravý sloupec
            %----------------------------------------
            \begin{column}{0.48\textwidth}
                \textbf{Konečný automat:}

                \begin{center}
                \begin{tikzpicture}[scale=0.95, every node/.style={font=\small}]
    \initialstate{q0}{(0,0)}{$q_0$}{(-1.0,0)}{west}
    \state{q1}{(2.0,0)}{$q_1$}
    \acceptingstate{q2}{(4.0,0)}{$q_2$}

    \transitionarrow[loop above]{q0}{above}{$b$}{q0}
    \transitionarrow[bend left=15]{q0}{above}{$a$}{q1}

    \transitionarrow[loop above]{q1}{above}{$a$}{q1}
    \transitionarrow[bend left=15]{q1}{above}{$b$}{q2}

    \transitionarrow[bend left=20]{q2}{below}{$a$}{q1}
    \transitionarrow[bend left=50]{q2}{above}{$b$}{q0}
\end{tikzpicture}

                \end{center}

                \begin{itemize}
                    \item \(q_0\): zatím nic podstatného
                    \item \(q_1\): poslední symbol je \(a\)
                    \item \(q_2\): poslední dva symboly jsou \(ab\)
                \end{itemize}
            \end{column}
        \end{columns}
    \end{frame}

    \section{Zdroje}

    \begin{frame}{Zdroje}
        \begin{itemize}
            \item PAPADIMITRIOU, Christos H. Computational complexity. Reading: Addison-Wesley, 1994. ISBN 0-201-53082-1.
            \item HOPCROFT, John E.; MOTWANI, Rajeev a ULLMAN, Jeffrey D. Introduction to automata theory, languages, and computation. 3rd ed. Boston: Pearson/Addison Wesley, c2007.
        \end{itemize}
    \end{frame}

\end{document}